\SourceLeaf{62}fortunate successes with three other new planets discovered since then.

Those methods are the original foundations of the work under review. In continuing to work on the subject, however, the author had the opportunity to make so many successive additions, changes and improvements to them that almost no trace of their original form remains. He therefore hopes that the astronomers who expressed a keen interest in these investigations from the outset and wished to see them published soon will have no reason to be dissatisfied with the delay. We leave it to competent readers of the work itself to judge how well founded that hope is and confine ourselves here to a very brief account of its plan and contents.

The problem mentioned repeatedly above was indeed the occasion for the present work and its chief aim, but in carrying it out the problem could occupy only part of it, and indeed the smaller part. It is easy to guess that an adequate solution could not be given without presupposing a host of truths about the motion of celestial bodies in conic sections that had yet to be developed. These new relations, already very interesting in themselves, were so numerous and so closely connected that they had to be presented together. The entire first part of the work is therefore devoted to them. Their arrangement follows a very simple plan. This first part is divided into four sections. The first and second contain relations concerning a single position of the celestial body, whether in its orbit or in space, while the third and fourth develop relations between two or more positions. Including the most important relations already known alongside the new ones throughout was hardly avoidable. Readers will find, however, that the known ones are usually developed in a form peculiar to the author. The limited space of these pages does not allow even an overview of all the subjects treated here. We confine ourselves to touching on a few of the more extensive investigations. In the first section there is a digression
\SourceLeaf{63}on the degree of accuracy attainable in numerical calculation with tables of all kinds. The same section also presents a new method for reducing motion in ellipses and hyperbolas that approach a parabola to motion in the latter. The second section gives a procedure for determining the geocentric positions of celestial bodies by means of three coordinates. As is well known, the author had already written a paper on it, and here the procedure is substantially improved. The same section also shows how to determine the differential changes in the geocentric position from the differential changes in the individual elements. The third section treats in exceptional detail the extremely important problem of determining the elements from two positions of a celestial body in its orbit. The second solution in particular provided an opportunity to develop a great many new relations. To show the practical usefulness of the methods more clearly and make them easier for the less experienced to study, the more important results are illustrated throughout with examples drawn from actual cases, most of them related to one another. Motion in a hyperbola, which can in fact also occur in comets, is treated almost as fully as motion in an ellipse or parabola for its analytical interest alone.

The investigations presented in the first part are, as it were, the material from which the second part constructs the solution to the great problem of determining the orbit of the celestial body itself by analyzing its geocentric appearances. This undertaking in fact involves several quite different demands. Finding a first approximation to a celestial body’s still entirely unknown orbit requires one kind of work. Refining elements already known approximately on the basis of a longer series of observations, perhaps spanning many years, calls for another. The first task requires greater skill, the other more labor. For the first one uses only as many observations as are needed, generally three. The first section of the second part is devoted to the problem of determining the orbit of a celestial body from three complete observations, in a sense the most important problem
\SourceLeaf{64}in the whole work. The solution itself cannot be abridged here. As an indication of how widely applicable it is, we mention only that in the three illustrative examples it is applied with equal success to observations 22 days apart, to ones 71 days apart and to ones 260 days apart. One distinctive feature is that, if desired, the small changes due to parallax and aberration can be taken into account \textit{at once}, without making the work significantly harder, so they do not require the calculation to be done twice.

Under certain circumstances, especially if the plane of the orbit is only slightly inclined to the ecliptic, it would be inadvisable to base the determination of the elements on three complete observations. Four observations must be used, of which only two need be complete. This gives rise to a new problem, to which the second section is devoted. An example taken from Vesta illustrates the method, with 162 days between the first and last observations. Here too, if desired, parallax and aberration can be taken into account \textit{at once}.

Refining the elements of a celestial body calls for the opposite procedure. The smallest possible number of observations gives way to as many as are available. The third section shows how to proceed. This was the place to set out the main points of a question of the highest importance for every application of mathematics to the physical world. Given the imperfection of our senses and instruments, observations and measurements always inevitably contain errors, however small. How are they best combined to establish results? The principles developed here, which the author had already been using for 14 years and had shared some considerable time ago with several of his friends in astronomy, lead to the method that \textsc{Legendre} also presented a few years ago in his \textit{Nouvelles méthodes pour la détermination des orbites des comètes} under the name \textit{Méthode des moindres carrés}. The justification that the author gives for the method is
\SourceLeaf{65}entirely his own. A fuller treatment from him is to be expected in due course.

To achieve \textit{lasting} agreement between the elements and the motions of a planet, one must go beyond purely elliptic motion and necessarily take into account the perturbations produced by the other planets. Calculating the perturbations themselves was of course outside the plan of a work devoted solely to motion according to \textsc{Kepler}’s laws. The fourth and final section does, however, at least indicate the main steps needed to improve a planetary orbit by taking account of perturbations already calculated, and so, as it were, to put the finishing touches to it.

We close this review with the wish that this work, the fruit of seven years spent on these investigations, may help win more friends for this most important and most beautiful part of theoretical astronomy and ease the work of astronomers themselves, whose tasks are growing in scope as the science steadily advances.
\par\bigskip\centerline{\rule{.24\textwidth}{.3pt}}
